\chapter{Logical Decision Problems and the Arithmetical Hierarchy}
\label{ch:decision-hierarchy}

\section{The classical decision problem}
\label{sec:decision-hierarchy-01}

Fix an effective coding of first-order formulas and finite derivations.

\begin{theorem}[Negative solution to the classical Entscheidungsproblem]
\label{thm:fol-validity-undecidable}
Under the fixed effective coding, no Turing machine decides the set of
codes of logically valid first-order sentences.
\textit{Status: cited theorem.}
\end{theorem}

Church and Turing supply the historical negative solution
\cite{Church1936,Turing1937}.  The theorem records only that historically
sourced general decision problem.

\begin{corollary}[Enumerability without decidability]
Under the fixed effective coding, the set of logically valid first-order
sentences is computably enumerable but not decidable.
\textit{Status: complete proof.}
\end{corollary}

\begin{proof}
Enumerate all finite derivations in the fixed effective first-order
calculus that have no open assumptions and whose conclusion is a sentence,
and output their conclusions.  Soundness for the empty premise set ensures
that every output is valid, and completeness ensures that every valid
sentence eventually appears as such a conclusion, so the
set is computably enumerable.  A decider for this set would solve the general
decision problem forbidden by \cref{thm:fol-validity-undecidable}.
\end{proof}

\begin{warning}[Restricted signatures]
The theorem concerns the general decision problem.  The monadic fragment
without function symbols is decidable, so undecidability may not be transferred
to an individual fixed signature or fragment without a theorem controlling
that restriction.
\end{warning}

\begin{theorem}[Trakhtenbrot]
For every finite relational signature containing a binary relation symbol, the
set of sentences having a finite model is computably enumerable but
undecidable, and the set of sentences true in every finite structure is not
computably enumerable.
\textit{Status: cited theorem.}
\end{theorem}

See \cite{Trakhtenbrot1950} for a complete proof.

The contrast with unrestricted validity is important: unrestricted
validities are computably enumerable by completeness of the fixed effective
first-order calculus, while finite validities admit no effective calculus
both sound and complete for truth in all finite
structures.

